\chapter{Equities and Options Connectivity}\label{ch:nw:equities-and-options-connectivity}

A firm that wants the whole U.S. options market must take a consolidated feed split across 96 multicast lines, sized by its
operator for bursts rather than averages, and send its quotes through ports whose number and message budgets it cannot
exceed. The feed's operator publishes its capacity in 10-millisecond peaks; the busiest millisecond of July 2026 held 0.35
million messages, about \qty{135}{\giga\bit\per\second} on the wire for one millisecond. Connectivity for options is capacity
planning before it is latency.

Book~1 introduced direct feeds, the securities information processors and sponsored access; Book~11 the mass quote, the
message throttle and the market access rule; Book~13 the feed handler. This chapter turns them into numbers of ports,
links and cores: the port types a venue sells, the bandwidth of the consolidated options feed, the risk layer that sits
between a sponsored firm and the venue, and the engineering of quote traffic. The plan is computed by \texttt{firm.portplan}
from the operator's and the venues' dated documents, with a simulated busy minute labelled as such.

\section{Gateway and port types}

\begin{definition}[Order-entry port]\label{def:nw:equities-and-options-connectivity:port}
An \emph{order-entry port}\index{order-entry port} is one logical connection, with its own session and credentials, through
which a member sends orders, quotes or cancels to a venue's gateway and receives their acknowledgements; venues sell
ports by type (order entry, bulk quoting, purge or mass cancel, drop copy) and charge and limit them per port.
\end{definition}

A venue's connectivity has two layers. The physical one is chapter~9's: a \omterm{def:nw:colocation-products-and-how-they-are-sold:mmr}{cross-connect}, a port speed, a \omterm{def:nw:colocation-products-and-how-they-are-sold:cabinet}{colocation
cabinet}. The logical one is the ports: sessions on the gateway, each with a type, a fee and rules. Options venues
separate quoting from ordinary order entry, because a market maker's quotes are the bulk of their inbound traffic.

The MIAX Express Interface shows the pattern. A Full Service port accepts every message, bulk quotes included; a Limited
Service port every message except bulk quotes; a Priority Mass Cancel port only mass cancels, processed ahead of the
firm's other traffic. Each matching environment (a ``cloud'') trades all options on a set of underlyings, and a firm gets
two Full Service and twelve Limited Service ports per cloud: to quote every underlying, it connects to every cloud. On
Nasdaq's Phlx the quoting interface is SQF, for quotes, immediate-or-cancel orders and auction responses; the exchange's own
filing says one port suffices to meet a market maker's quoting obligations, and its fee schedule caps a market maker at 250
of them.

\begin{dated}{2026-09}{Quoting ports on U.S. options exchanges, from the venues' filings and specifications}\label{dat:nw:equities-and-options-connectivity:ports}
\begin{itemize}
\item Phlx: SQF port \$1\,185 a port a month from January 2026, 10\% to 50\% off by market-maker volume share; at most 250
  SQF ports per market maker. BOX: \$1\,000 a month for all ports. MIAX: MEI fees by classes quoted, up to \$20\,500 a month
  for over 100 classes. NYSE Arca and NYSE American: \$510 a port for ports 1 to 40, \$170 from port 41. Cboe EDGX Options:
  \$750 a logical port (all as compared in Phlx's September 2025 filing).
\item MIAX MEI: 2 Full Service and 12 Limited Service ports per cloud; a Simple Bulk Quote message carries up to 50
  single-sided quotes, 51 bytes of header and 15 bytes a quote; the specification calls MEI ``a synchronous quoting
  interface''.
\end{itemize}
\end{dated}

Twenty-four ports cost \$28\,440 a month on Phlx, \$18\,000 on EDGX, \$12\,240 on Arca or American, \$1\,000 on BOX, and on
MIAX the fee depends on the classes, not the ports (\$20\,500 for all of them). The drop copy and the purge ports come on
top. Fees differ by a factor of twenty-eight for the same count of sessions: the count is a design choice, and the fee
schedule is one of its inputs.

\section{Consolidated options-feed bandwidth}

\begin{definition}[Feed partition]\label{def:nw:equities-and-options-connectivity:partition}
A \emph{feed partition}\index{feed partition} is one of the parallel streams (lines, channels, \omterm{def:nw:networking-for-trading:group}{multicast groups}) into which a
feed's operator splits its messages by a published rule on the instrument, so that each stream fits a link and a handler
and a subscriber can take only the partitions it needs.
\end{definition}

The consolidated options feed has 96 lines since its expansion, and since 1 June 2026 a published symbol distribution
spreads the load over them: by symbol range, and within the busiest symbols by calls and puts and by odd and even expiration
months or days, ``for optimal symbol balancing and line capacity utilization''. Its capacity notice is written for burst
planning (\cref{dat:nw:equities-and-options-connectivity:opra}).

\begin{dated}{2026-09}{The consolidated options feed's capacity, from its operator's notices}\label{dat:nw:equities-and-options-connectivity:opra}
The projection effective July 2026 gives for one stream 13.575 million messages, 4.403 gigabits and 1.564 million packets
in the busiest 100 milliseconds; 1.562 million messages and 0.501 gigabits in the busiest 10; 311 billion messages a day;
and at most 625\,000 messages per 100 milliseconds on one line. Both redundant streams double the bandwidth; add 10\% for
retransmissions. The median latency through the processor is under 18 microseconds. Ninety-six lines, with a new symbol
distribution from 1 June 2026.
\end{dated}

Read as rates, the notice asks for \qty{44.0}{\giga\bit\per\second} over the busiest 100 milliseconds and
\qty{50.1}{\giga\bit\per\second} over the busiest 10. Its gigabits divided by its packets give 352 bytes a packet and 8.68
messages a packet; treating those bytes as UDP payload and adding Ethernet, IP, UDP and the line's own overhead with
\texttt{firm.netsim} gives 418 bytes on the wire a packet, \qty{48.0}{bytes} a message. On the wire the 10-millisecond peak
is \qty{59.4}{\giga\bit\per\second} a stream, and with both streams and the retransmission allowance
\qty{130.6}{\giga\bit\per\second}.

A switch port does not average over 10 milliseconds. To plan links, the chapter simulates a busy minute in 1-millisecond
bins: a slowly varying level (a log-AR(1) with a persistence of a second) times rare spikes (one bin in a thousand) and a
little noise. Two parameters, the level's spread and the spikes' size, are fitted to two published ratios: the busiest
millisecond of July 2026 against the busiest 10 milliseconds, 3.93, and the projection's busiest 10 milliseconds against its
busiest 100, 1.151. The fit gives 3.90 and 1.152; scaled to the projection's 10-millisecond peak, the simulated minute's
busiest 100 milliseconds hold 13.56 million messages against the notice's 13.575, a check the fit was not asked to pass. Each
bin's messages are then shared over the 96 lines with lognormal noise around equal shares (spread 0.5, an assumption).

\omcode{code/firm/portplan/firm_portplan.py}{74}{82}{The simulated busy minute: a slow level times rare spikes, in 1-ms bins.}\label{lst:nw:equities-and-options-connectivity:burst}

\begin{omfigure}
\begin{tikzpicture}
\begin{axis}[width=11.5cm,height=5.6cm,xmin=-300,xmax=300,ymode=log,ymin=8,ymax=400,
  xlabel={time from the busiest millisecond (ms)},ylabel={one stream on the wire (Gb/s, log)},
  tick label style={font=\footnotesize},label style={font=\small},grid=major,grid style={gray!25},
  legend style={font=\footnotesize,at={(0.5,-0.3)},anchor=north,legend columns=3},table/col sep=comma]
\addplot[omThm!70,thin] table[x=t_ms,y=gbps_1ms]{figdata/networks/22-equities-and-options-connectivity/burst.csv};
\addplot[omDef,thick] table[x=t_ms,y=gbps_10ms]{figdata/networks/22-equities-and-options-connectivity/burst.csv};
\addplot[omMeth,dashed] coordinates {(-300,10) (300,10)};
\legend{1-ms bins,10-ms average,one 10 Gb/s link}
\end{axis}
\end{tikzpicture}

\omcaption{Simulation: six hundred milliseconds of the busy minute around its busiest millisecond, one stream, on the wire. The
busiest millisecond runs at \qty{234}{\giga\bit\per\second}, the 10-millisecond average around it below
\qty{40}{\giga\bit\per\second}. Data: \texttt{fig\_options.py}, \texttt{nw\_options.LINES\_GBPS}.}\label{fig:nw:equities-and-options-connectivity:burst}
\end{omfigure}

The lines are then packed onto links, whole and in order, and a link is enough at a planning percentile when that
percentile of its 1-millisecond rate fits \qty{10}{\giga\bit\per\second} less the 10\% allowance.

\omcode{code/firm/portplan/firm_portplan.py}{113}{121}{Links for one stream: the fewest groups of whole lines whose planning percentile fits a link.}\label{lst:nw:equities-and-options-connectivity:links}

\begin{proposition}[Links against the percentile]\label{prop:nw:equities-and-options-connectivity:links}
The number of links a feed needs never falls when the planning percentile rises, and at the 100th percentile of 1-millisecond
bins it is at least the feed's busiest millisecond divided by a link's usable capacity.
\end{proposition}

\begin{proof}
A higher percentile of each group's rate is at least the lower one, so every packing that fails at the lower percentile fails
at the higher. At the maximum, each link carries at most its capacity in every bin, so the links together carry at most their
number times the capacity in the busiest bin, which must hold the whole feed's busiest millisecond.
\end{proof}

\begin{omfigure}
\begin{tikzpicture}
\begin{axis}[ybar,width=11cm,height=5.8cm,bar width=9pt,ymin=0,ymax=200,xtick=data,
  xticklabels from table={figdata/networks/22-equities-and-options-connectivity/links.csv}{pct},
  xlabel={planning percentile of 1-ms bins},ylabel={links, both streams},
  nodes near coords,every node near coord/.append style={font=\scriptsize},
  tick label style={font=\footnotesize},label style={font=\small},grid=major,grid style={gray!25},
  legend pos=north west,legend style={font=\footnotesize},table/col sep=comma]
\addplot[fill=omThm!45,draw=omThm] table[x=x,y=links10]{figdata/networks/22-equities-and-options-connectivity/links.csv};
\addplot[fill=omDef!55,draw=omDef] table[x=x,y=links25]{figdata/networks/22-equities-and-options-connectivity/links.csv};
\legend{10 Gb/s links,25 Gb/s links}
\end{axis}
\end{tikzpicture}

\omcaption{Simulation: links needed for both streams of the consolidated options feed, 96 lines packed whole, against the
percentile of 1-millisecond bins the plan must carry. Median: 4; 99th: 12; 99.99th: 44; the busiest millisecond: 178 links
of \qty{10}{\giga\bit\per\second} or 32 of \qty{25}{\giga\bit\per\second}. Data: \texttt{fig\_options.py},
\texttt{nw\_options.links\_table()}.}\label{fig:nw:equities-and-options-connectivity:links}
\end{omfigure}

The answer depends on the percentile more than on anything else. At the median, four \qty{10}{\giga\bit\per\second} links carry
both streams; at the 99th percentile twelve; at the 99.99th, which leaves six milliseconds of the minute over budget,
forty-four. The proposition's bound for the busiest millisecond is 26 links a stream (\qty{234}{\giga\bit\per\second} over
\qty{9.09}{\giga\bit\per\second}); packing whole lines needs 89, 178 for both streams, nearly one link a line, because a single
line's busiest millisecond reaches \qty{6.8}{\giga\bit\per\second} and the lines' bursts do not fall in the same millisecond.
With \qty{25}{\giga\bit\per\second} links, 32. What the plan does not carry at the link, a
switch buffer or a drop absorbs (chapter~1), so the percentile is a choice of how often the firm accepts a queue or a gap.
The line limit is not the constraint: the busiest line's busiest 100 milliseconds hold 171\,500 messages, a quarter of the
625\,000 the notice allows.

The handlers are the second budget. At 10 million messages a second a core (an assumption: \qty{100}{\nano\second} a message, well above
Book~13's decoding cost of a few nanoseconds, to leave room for building the book), the 10-millisecond peak rate of 156.2 million messages a second
needs 16 cores a stream. They cannot absorb the busiest millisecond, 609\,500 messages, without a queue.

\omcode{code/networks/22-equities-and-options-connectivity/python/nw_options.py}{37}{44}{The queue in front of the handler cores, as a fluid over the 1-ms bins.}\label{lst:nw:equities-and-options-connectivity:backlog}

With 16 cores the queue reaches 449\,500 messages and takes \qty{2.8}{\milli\second} to clear; with 32, \qty{0.9}{\milli\second};
with 48, \qty{0.27}{\milli\second}; 64 cores never queue. Cores bought for the 10-millisecond peak turn the busiest millisecond
into milliseconds of stale books, exactly when the market moves.

\section{Risk-check layers under sponsored access}

\begin{definition}[Sponsored-access risk layer]\label{def:nw:equities-and-options-connectivity:risk}
A \emph{sponsored-access risk layer}\index{sponsored-access risk layer} is the sponsoring broker's pre-trade control, in software
or in a hardware gateway, through which every order and quote of a sponsored firm passes on its way to the venue, and
which adds its processing latency to every message.
\end{definition}

Under the market access rule (Book~11), a broker that gives a firm access to a venue under its membership must check that
firm's messages before they reach the venue. The check sits on the path, and on a quoting interface it sits on every round
trip. Vendors publish their gateways' latency.

\begin{dated}{2026-09}{Published latencies of hardware risk gateways}\label{dat:nw:equities-and-options-connectivity:risk}
Magmio's FPGA Risk Check Gateway: pre-trade risk and compliance checks ``as low as 200 nanoseconds''. Algo-Logic's FPGA
Pre-Trade Risk Check: checks ``well under one microsecond'', for Rule 15c3-5 compliance.
\end{dated}

No broker publishes the latency of its software layer; the chapter assumes \qty{5}{\micro\second} each way. The cost of a
risk layer is not only the added microseconds on one message: on a synchronous interface, it lowers what a port can carry.

\section{Quote-traffic engineering}

A market maker's quoting load is its classes times their series times the update rate times two sides. Bulk messages carry
up to 50 single-sided quotes, so messages are fifty times fewer than quotes. On a synchronous interface, read here as one
bulk message in flight per port until its acknowledgement (an assumption to confirm with the venue), a port sends one
message per round trip.

\omcode{code/firm/portplan/firm_portplan.py}{132}{139}{Ports: one message per round trip on each, and at least one per matching engine.}\label{lst:nw:equities-and-options-connectivity:ports}

\begin{proposition}[Quote refresh time]\label{prop:nw:equities-and-options-connectivity:refresh}
On a synchronous quoting interface with $k$ ports per matching engine, $b$ quotes per message and a round trip $\tau$ (risk
layer included), refreshing $Q$ quotes spread evenly over $E$ engines takes $Q\tau/(Ekb)$: linear in the round trip and in
the inverse of the ports.
\end{proposition}

\begin{proof}
Each engine receives $Q/E$ quotes in $Q/(Eb)$ messages; its $k$ ports send $k$ messages per round trip in parallel, so they need
$Q/(Ekb)$ round trips.
\end{proof}

\begin{omfigure}
\begin{tikzpicture}
\begin{axis}[width=11cm,height=5.6cm,xmin=0,xmax=10,ymin=0,ymax=42,
  xlabel={risk layer, one way (\unit{\micro\second})},ylabel={time to refresh every quote (ms)},
  tick label style={font=\footnotesize},label style={font=\small},grid=major,grid style={gray!25},
  legend pos=north west,legend style={font=\footnotesize},table/col sep=comma]
\addplot[omThm,thick] table[x=risk_us,y=one_port]{figdata/networks/22-equities-and-options-connectivity/refresh.csv};
\addplot[omDef,thick] table[x=risk_us,y=two_ports]{figdata/networks/22-equities-and-options-connectivity/refresh.csv};
\legend{one quoting port per engine,two quoting ports per engine}
\end{axis}
\end{tikzpicture}

\omcaption{The time to refresh all 600\,000 quotes of 1\,200 classes of 250 series over twelve matching engines, 50 quotes a
message, one message in flight per port, a \qty{20}{\micro\second} round trip before the risk layer (assumptions). Two ports and
no risk layer: \qty{10}{\milli\second}; a \qty{5}{\micro\second} software layer: 15. Data: \texttt{fig\_options.py},
\texttt{nw\_options.refresh\_curve()}.}\label{fig:nw:equities-and-options-connectivity:refresh}
\end{omfigure}

The chapter's firm quotes 1\,200 classes of 250 series on average, four updates a second a side: 2.4 million quotes a second,
48\,000 bulk messages, \qty{0.31}{\giga\bit\per\second} of order entry. With a \qty{20}{\micro\second} round trip a port carries
50\,000 messages a second, so throughput alone asks for one port; the venue's twelve engines ask for twelve, and the two quoting
ports per engine the venue allocates give twenty-four. Order entry is bound by messages and engines, not bandwidth; the feed is
bound by bandwidth.

The binding case is the market-wide move, when every quote must change at once: 600\,000 quotes. Two ports an engine refresh them
in \qty{10.0}{\milli\second}; behind the FPGA gateway, \qty{10.2}{\milli\second}; behind the assumed software layer,
\qty{15.0}{\milli\second}; with one port an engine and no layer, \qty{20.0}{\milli\second}. Quote-traffic engineering is the
art of not needing that full refresh: single-sided quotes let the firm send the important sides first (the specification's own
point), purge ports pull quotes faster than updating them, and the order in which series are refreshed is a risk decision.

\begin{method}[Sizing an options desk's connectivity]\label{met:nw:equities-and-options-connectivity:size}
\begin{enumerate}
\item Read the feed's latest capacity notice and metrics; convert them to wire bits with the framing overhead; double for both
  streams and add the retransmission allowance.
\item Choose a planning percentile at the millisecond scale, and pack whole partitions onto links against it; decide what the
  switch buffers absorb above it.
\item Size handler cores on the 10-millisecond peak and check the queue of the busiest millisecond; add cores until the delay is
  acceptable.
\item From the quoting load, compute messages, ports per engine and the full-refresh time, with the risk layer in the round
  trip.
\item Price the ports on each venue's schedule; add drop copies and purge ports; re-run at each new capacity notice.
\end{enumerate}
\end{method}

\begin{table}[ht]
\centering\small
\begin{tabular}{@{}llr@{}}
\toprule
Budget & Quantity & Plan \\
\midrule
Feed, one stream & 10-ms peak on the wire & \qty{59.4}{\giga\bit\per\second} \\
Feed, both streams & with 10\% retransmissions & \qty{130.6}{\giga\bit\per\second} \\
Links & 99.99th percentile, 10 Gb/s & 44 \\
Handler cores & 10-ms peak, one stream & 16 \\
Busiest millisecond & queue with 16 cores & \qty{2.8}{\milli\second} \\
Quoting & quotes and messages a second & 2.4 million; 48\,000 \\
Quoting ports & two per engine, twelve engines & 24 \\
Full refresh & two ports, software risk layer & \qty{15.0}{\milli\second} \\
\bottomrule
\end{tabular}
\caption{The capacity plan of the chapter's options desk (feed from the July 2026 projection and a labelled simulation; quoting
from stated assumptions). Data: \texttt{nw\_options.feed\_plan()}, \texttt{quote\_plan()}.}\label{tab:nw:equities-and-options-connectivity:plan}
\end{table}

\begin{omfigure}
\begin{tikzpicture}[font=\small,
  box/.style={draw=omDef,fill=omDef!10,rounded corners,align=center,minimum height=9mm,inner sep=3pt},
  ven/.style={draw=omThm,fill=omThm!10,rounded corners,align=center,minimum height=9mm,inner sep=3pt}]
\node[ven] (opra) at (0,2.2) {consolidated feed\\96 lines, 2 streams};
\node[box] (sw) at (3.9,2.2) {switch\\links per percentile};
\node[box] (fh) at (7.6,2.2) {handler cores\\and queue};
\node[box] (st) at (10.6,1.1) {strategy};
\node[box] (rk) at (7.6,0) {risk layer\\(broker)};
\node[box] (pt) at (3.9,0) {quoting ports\\per engine};
\node[ven] (me) at (0,0) {matching\\engines};
\draw[->,omThm] (opra) -- (sw); \draw[->,omThm] (sw) -- (fh); \draw[->,omThm] (fh) -| (st);
\draw[->,omDef] (st) |- (rk); \draw[->,omDef] (rk) -- (pt); \draw[->,omDef] (pt) -- (me);
\draw[->,omDef,dashed] (me.south) -- ++(0,-0.5) -| node[pos=0.25,below,font=\footnotesize] {acknowledgements: one round trip per bulk message} (rk.south);
\end{tikzpicture}

\omcaption{The two budgets of an options desk: the feed side is bound by bandwidth and handler cores, the quoting side by
ports, engines and round trips, with the broker's risk layer on every one. Schematic.}\label{fig:nw:equities-and-options-connectivity:desk}
\end{omfigure}

\section{Tutorial: a capacity table for an options desk}\label{tut:nw:equities-and-options-connectivity:tut}

\begin{tutorial}
\textbf{Goal.} From the feed's capacity notice and a quoting load, compute links, cores and ports.
\textbf{End state:} \cref{tab:nw:equities-and-options-connectivity:plan} and \cref{fig:nw:equities-and-options-connectivity:links}.
\begin{tutsteps}
\item \textbf{Notice.} \texttt{firm\_portplan.projection} reads the dated rows for July 2026; \texttt{wire\_bytes\_per\_msg} adds the
  framing with \texttt{firm.netsim}.
\item \textbf{Busy minute.} \texttt{fit\_burst} fits the burst model to the two published ratios; \texttt{burst}
  (\cref{lst:nw:equities-and-options-connectivity:burst}) and \texttt{line\_split} give each line's 1-ms rates.
\item \textbf{Links and cores.} \texttt{links\_needed} (\cref{lst:nw:equities-and-options-connectivity:links}) for each
  percentile; \texttt{nw\_options.backlog} (\cref{lst:nw:equities-and-options-connectivity:backlog}) for the cores.
\item \textbf{Ports.} \texttt{ports\_needed} (\cref{lst:nw:equities-and-options-connectivity:ports}) and
  \texttt{nw\_options.quote\_plan} with each risk layer.
\end{tutsteps}
\textbf{What to change next.} Replace the equal line shares with the published symbol distribution and a model of each
symbol's traffic; add the switch buffer with \texttt{firm.netsim.egress} and count drops at each percentile.
\end{tutorial}

\section{Build: the port and bandwidth planner}\label{bld:nw:equities-and-options-connectivity:portplan}

\begin{build}
\bfield{Purpose} Links, cores and ports for a desk that takes a partitioned feed and quotes through a venue's ports: the
equities and options rows of chapter~29's plan.
\bfield{Interface} \texttt{firm\_portplan}: \texttt{load\_capacity}, \texttt{projection}, \texttt{wire\_bytes\_per\_msg},
\texttt{burst}, \texttt{window\_peak}, \texttt{ratios}, \texttt{fit\_burst}, \texttt{line\_split}, \texttt{links\_needed},
\texttt{cores\_needed}, \texttt{quotes\_per\_s}, \texttt{port\_msgs\_per\_s}, \texttt{ports\_needed}, \texttt{bulk\_bytes}.
\bfield{Rules} Published capacities and fees are dated rows with sources; the busy minute is a labelled simulation whose free
parameters are fitted to published ratios; every other parameter is a stated assumption.
\bfield{Acceptance tests} \texttt{code/firm/portplan/tests/}: the capacity rows, flat traffic giving ratios of one, equal
lines packing by hand, and the port arithmetic.
\bfield{Stretch} Line shares from the published symbol distribution; drops through \texttt{firm.netsim}'s \omterm{def:nw:networking-for-trading:egress}{egress queue}; per-venue
port rules.
\end{build}

\begin{omsources}
\item SIAC, OPRA capacity projections (September 2025) and 96-line symbol distribution notice (May 2026); OPRA key operating
  metrics (August 2026).
\item SEC Release No.~34-103889 and Federal Register 90 FR 60788 (Phlx SQF port fees); MIAX Options, MEI Interface
  Specification, version 2.10a.
\item Magmio and Algo-Logic product pages for FPGA risk gateways.
\end{omsources}

\section{Exercises}

\begin{exercise}[$\star$]\label{exo:nw:equities-and-options-connectivity:1}
From the July 2026 projection, what are the busiest 100 milliseconds and the busiest 10 milliseconds as rates in gigabits a
second, for one stream as the notice counts it?
\end{exercise}

\begin{exercise}[$\star$]\label{exo:nw:equities-and-options-connectivity:2}
How many bulk messages a second does a firm send for 2.4 million single-sided quotes a second on MEI?
\end{exercise}

\begin{exercise}[$\star$]\label{exo:nw:equities-and-options-connectivity:3}
Why does a firm quoting on MIAX need ports on every cloud, even with little load on some of them?
\end{exercise}

\begin{exercise}[$\star\star$]\label{exo:nw:equities-and-options-connectivity:4}
Using \cref{prop:nw:equities-and-options-connectivity:refresh}, how long does a full refresh take with two ports an engine behind the
FPGA gateway, and what would a third port an engine give?
\end{exercise}

\begin{exercise}[$\star\star$]\label{exo:nw:equities-and-options-connectivity:5}
Why is the 99.99th percentile a demanding planning target at the millisecond scale, and what happens above it?
\end{exercise}

\begin{exercise}[$\star\star$]\label{exo:nw:equities-and-options-connectivity:6}
What does twenty-four quoting ports cost a month on each venue in \cref{dat:nw:equities-and-options-connectivity:ports}?
\end{exercise}

\begin{exercise}[$\star\star\star$]\label{exo:nw:equities-and-options-connectivity:7}
\emph{Coding.} With \texttt{nw\_options.backlog}, how many handler cores keep the busiest millisecond's queue under
\qty{0.5}{\milli\second}?
\end{exercise}

\begin{exercise}[$\star\star\star$]\label{exo:nw:equities-and-options-connectivity:8}
\emph{Find the flaw.} ``The notice says \qty{50}{\giga\bit\per\second} at the 10-millisecond peak, so six
\qty{10}{\giga\bit\per\second} links per stream are plenty.''
\end{exercise}

\section{Problem: Quoting Every Series}

\begin{problem}[{Weekend problem --- ports and links for a U.S. options market maker}]\label{pb:nw:equities-and-options-connectivity:1}
A firm will quote 1\,200 option classes of 250 series on average, four updates a second a side, on a venue with twelve matching
engines and synchronous bulk quoting, through a sponsoring broker; it takes the consolidated feed in full, both streams. Use the
chapter's planner.

\medskip\noindent\textbf{Part I --- The feed.}
\begin{enumerate}
\item What bandwidth does the July 2026 projection ask for, as the notice counts it and on the wire, one stream and both with the
  allowance?
\item How many bytes a message does the chapter use on the wire, and from what?
\item How many 10 Gb/s links carry both streams at the median, 99th and 99.99th percentiles of 1-millisecond bins?
\item Why does carrying the busiest millisecond need 178?
\end{enumerate}

\medskip\noindent\textbf{Part II --- The handlers.}
\begin{enumerate}[resume]
\item How many cores does the 10-millisecond peak need at 10 million messages a second a core?
\item What queue does the busiest millisecond build with them, and how long does it take to clear?
\item How many cores remove the queue entirely?
\item Which check did the burst model pass without being fitted to it?
\end{enumerate}

\medskip\noindent\textbf{Part III --- The quotes.}
\begin{enumerate}[resume]
\item How many quotes, bulk messages and gigabits a second does the firm send?
\item How many ports does throughput need, and how many does the venue's structure impose?
\item How long does a full refresh take with two ports an engine, without a risk layer, behind the FPGA gateway and behind the
  software layer?
\item What does the full refresh time mean for the firm's risk after a market-wide move?
\end{enumerate}

\medskip\noindent\textbf{Part IV --- The verdict.}
\begin{enumerate}[resume]
\item State the \emph{named result}: the ports and 10 Gb/s links needed to quote the firm's classes at its update rate and to take
  the consolidated feed in full.
\item Which inputs are published and which assumed?
\item What would 25 Gb/s links change?
\item Which of the firm's decisions does the risk layer affect most?
\item How should the firm order its quote updates after a move?
\item What must be re-run when the next capacity notice arrives?
\item Where do these numbers go in chapter~29's plan?
\item In one sentence: why is options connectivity capacity planning before it is latency?
\end{enumerate}
\end{problem}

\section{Interview questions}

\begin{interviewq}[$\star$ \iqroles{developer}]\label{iq:nw:equities-and-options-connectivity:1}
Why does a feed operator publish its capacity per 10 milliseconds rather than per second?
\end{interviewq}

\begin{interviewq}[$\star\star$ \iqroles{developer}]\label{iq:nw:equities-and-options-connectivity:2}
How would you decide how many links and servers to buy for the full U.S. options feed?
\end{interviewq}

\begin{interviewq}[$\star\star$ \iqroles{developer, researcher}]\label{iq:nw:equities-and-options-connectivity:3}
What limits how fast a market maker can update all its quotes, and how would you measure it?
\end{interviewq}

\begin{interviewq}[$\star\star$ \iqroles{developer}]\label{iq:nw:equities-and-options-connectivity:4}
A sponsoring broker offers a software risk layer or a hardware gateway. What would you ask, and how would you compare them?
\end{interviewq}

\begin{interviewq}[$\star\star$ \iqroles{trader}]\label{iq:nw:equities-and-options-connectivity:5}
After a large move, which quotes should be updated first, and why?
\end{interviewq}

\begin{interviewq}[$\star\star\star$ \iqroles{developer, researcher}]\label{iq:nw:equities-and-options-connectivity:6}
Design the connectivity of a new options market-making desk in one U.S. data centre: feed, handlers, ports, risk layer and
monitoring.
\end{interviewq}
